% SPDX-License-Identifier: Apache-2.0
%
% The project's dynamics: how the work actually went. Built by
% //docs:dynamics. One column, because the chart is wide.
\documentclass[journal,onecolumn,9pt]{IEEEtran}

\newcommand{\listingsize}{\normalsize}
% SPDX-License-Identifier: Apache-2.0
%
% The house preamble, shared by every document in this directory. Loaded
% after \documentclass, before \title. Keeping it in one file is what
% stops two articles drifting apart in font, listing style or figure
% style.
% Computer Modern. IEEEtran selects Times (ptm) by default, so the face has
% to be chosen explicitly.
%
% Latin Modern is the Computer Modern variety used here, and the choice is
% forced rather than aesthetic. Under T1 encoding, plain `cmr` has no Type 1
% outlines unless cm-super is installed, and it is not in the pinned
% distribution, so pdflatex falls back to EC bitmaps and every glyph in the
% PDF becomes a Type 3 font. That was measured: the first build produced a
% document with 10 Type 3 fonts and no Type 1. Latin Modern is Computer
% Modern redrawn with a full T1 repertoire and real .pfb outlines, so the
% same design comes out scalable.
\usepackage[T1]{fontenc}
\usepackage[utf8]{inputenc}
\usepackage{lmodern}
% microtype is not in the pinned TeX distribution. emergencystretch alone
% keeps the two-column measure from producing overfull lines.
\setlength{\emergencystretch}{3em}

\usepackage{amsmath}
\usepackage{amssymb}
\usepackage{amsthm}
\usepackage{booktabs}
\usepackage{array}
% enumitem is not in the pinned TeX distribution; the list spacing below
% does what \setlist would have done.
\usepackage{float}
\usepackage{xcolor}
\usepackage{listings}
\usepackage{tikz}
\usetikzlibrary{arrows.meta,positioning,fit,backgrounds,calc}
\usepackage[hidelinks,bookmarks=true]{hyperref}

% Numbered, definition-style box for a rule the reader is meant to apply.
\theoremstyle{definition}
\newtheorem{rulebox}{Rule}
\newtheorem{example}{Example}

% Unnumbered asides.
\theoremstyle{remark}
\newtheorem*{tip}{Tip}
\newtheorem*{pitfall}{Pitfall}

% Tighter lists, in place of enumitem's \setlist.
\setlength{\topsep}{2pt}
\setlength{\itemsep}{1pt}
\setlength{\parsep}{0pt}

% Code in the running text. The typewriter face under T1 ligates `<<`
% and `>>` into guillemets, so a shift printed as \code{<<} came out as
% a French quotation mark (issue 571). microtype's \DisableLigatures is
% not in the pinned distribution, but the primitive it rests on is
% pdfTeX's own: \pdfnoligatures strips every ligature from the font it
% is given, and the current font inside \texttt is the typewriter face
% at the size in use, so each size is stripped the first time \code
% sets it. Code wants no ligature at all: two quotes are two quotes.
\newcommand{\code}[1]{\texttt{\ifdefined\pdfnoligatures\pdfnoligatures\font\fi #1}}
\newcommand{\term}[1]{\emph{#1}}

% Syntax highlighting. Four roles, distinguishable in colour and still
% distinguishable in greyscale, because an IEEE article gets printed: the
% keyword is the only bold role, the comment is the only italic one, and the
% two remaining roles differ in lightness as well as in hue.
\definecolor{kwblue}{RGB}{20,60,150}
\definecolor{cmtgreen}{RGB}{90,120,90}
\definecolor{strred}{RGB}{150,50,40}
\definecolor{numgray}{gray}{0.45}
\definecolor{framegray}{gray}{0.70}
\definecolor{bgshade}{gray}{0.97}
% A darker fill for figure cells. The listing background has to stay very
% light to keep code readable; a figure cell has to be clearly shaded or the
% distinction it draws is invisible in print.
\definecolor{cellshade}{gray}{0.90}

% The listing size is the document's choice, because it depends on the
% column: a two-column article sets `\listingsize` to 6pt and keeps its
% listings within 70 columns, or to 5.6pt when it includes source that
% rustfmt sets at 80, which is what fits a column's frame (issue 1061);
% a one-column document keeps the body size and includes source files
% at 80. `//docs:frame_test` checks every listing line against its frame. Lines are never broken; a line that does not fit
% overflows the frame, which is visible, rather than wrapping, which is
% not.
\providecommand{\listingsize}{\scriptsize}

% `'` is deliberately not a string delimiter: the language uses it for loop
% labels, as Rust does, so treating it as a quote would swallow the rest of
% the line.
\lstdefinestyle{txhdl}{
  basicstyle=\ttfamily\listingsize,
  keywordstyle=\bfseries\color{kwblue},
  commentstyle=\itshape\color{cmtgreen},
  stringstyle=\color{strred},
  numberstyle=\tiny\color{numgray},
  morekeywords={mod,use,pub,const,type,struct,enum,trait,impl,fn,async,let,mut,
    match,if,else,loop,while,for,in,break,continue,return,as,await,
    module,bus,channel,transaction,protocol,pipeline,flow,seq,config,
    port,stage,pipe,instance,connect,bind,tag,domain,blackbox,foreign,
    property,role,signal,handler,true,false,
    sequence,interface,modport,implements,package,import,var,wire,func,proc,
    case,default,next,fork,branch,join,state,fsm},
  morecomment=[l]{//},
  morestring=[b]",
  showstringspaces=false,
  columns=fullflexible,
  keepspaces=true,
  breaklines=false,
  % Framed, so a listing is bounded rather than floating in the column.
  frame=single,
  rulecolor=\color{framegray},
  framesep=3pt,
  backgroundcolor=\color{bgshade},
  xleftmargin=3pt,
  xrightmargin=3pt,
  aboveskip=6pt,
  belowskip=6pt,
  % Every listing is numbered and captioned, so it can be referred to by
  % number. `caption` without `float` numbers the listing and leaves it
  % where it was written, which is what a listing in running prose needs.
  captionpos=b,
  abovecaptionskip=2pt,
  belowcaptionskip=6pt,
}

% Figures drawn as text. Framed the same way, and with the highlighting
% switched off, because a box-and-arrow diagram is not code and half its
% words turning blue reads as an accident. Setting a style adds keys rather
% than clearing the ones already set, so the three roles are neutralised by
% hand rather than by leaving them out. Program output is printed in this
% style too, and a netlist's assignment can run past the frame, so a line
% that does not fit is broken and the rest indented; source never is.
\lstdefinestyle{ascii}{
  basicstyle=\ttfamily\listingsize,
  keywordstyle=\ttfamily,
  commentstyle=\ttfamily,
  stringstyle=\ttfamily,
  columns=fullflexible,
  keepspaces=true,
  breaklines=true,
  breakindent=2em,
  frame=single,
  rulecolor=\color{framegray},
  framesep=3pt,
  backgroundcolor=\color{bgshade},
  xleftmargin=3pt,
  xrightmargin=3pt,
  aboveskip=6pt,
  belowskip=6pt,
}
\lstset{style=txhdl}

% House TikZ styles. `shorten` lives on the arrow styles rather than on each
% arrow, so every arrow in the document gets the gap, including ones added
% later. TikZ clips a path at a node's border, so without it an arrowhead
% butts against the box with no gap at all. Keep `node distance` well above
% twice the shorten, or the arrow shrinks to nothing.
\tikzset{
  box/.style={draw,rounded corners=2pt,align=center,inner sep=4pt,
              font=\footnotesize},
  boxf/.style={box,fill=bgshade},
  lbl/.style={font=\scriptsize\itshape,align=center},
  fl/.style={-{Stealth[length=4pt]},shorten >=2.5pt,shorten <=2.5pt},
  fld/.style={fl,dashed},
}


% The contents list: the class gives a section's number 2.75em, which
% a Roman numeral past thirty overruns onto the title, so the box is
% wider here, and a subsection's entry starts where a title does.
\makeatletter
\def\l@section#1#2{\addpenalty{\@secpenalty}\addvspace{1.0em plus 1pt}%
    \@tempdima 5.75em \begingroup \parindent \z@ \rightskip \@pnumwidth%
    \parfillskip-\@pnumwidth {\bfseries\leavevmode #1}\hfil\hbox to\@pnumwidth{\hss #2}\par%
    \endgroup}
\def\l@subsection{\@dottedtocline{2}{5.75em}{5em}}
\def\l@subsubsection{\@dottedtocline{3}{10.75em}{5.5em}}
\makeatother


\InputIfFileExists{buildstamp.tex}{}{}
\input{timeline_counts}
\providecommand{\buildstamp}{Draft build (outside the Bazel flow).}

\title{The Dynamics of a Project}

\author{Filip Filmar and Dragiša Janković%
\thanks{Both authors are independent. Filip Filmar is the
corresponding author, at f@hdlfactory.com; Dragiša Janković is
reached at linkedin.com/in/dragisa-jankovic.}%
\thanks{This document is about how TxHDL was built rather than what
it is. What it is, is in the paper~\cite{paper}; the
cover~\cite{cover} lists every document. The chart and every number
here are measured from the repository's own history, not typed in:
\code{//tools/timeline} measured it on September 16, 2026, run by hand
since a Bazel action has no repository to read, and
\code{//tools/gantt} draws it at build time.}%
\thanks{The authors used a large language model, Claude, as an
assistant in exploring the concepts and in writing and constructing
the documents and the programs; every commit of the repository records
this attribution and includes the exact prompts that guided it.}%
\thanks{\buildstamp}}

\markboth{The Dynamics of a Project}{Filmar and Janković: Dynamics}

\begin{document}
\maketitle

\begin{abstract}
A hardware description language, its runtime, its lowering to two
netlist languages, a RISC-V core that runs on an FPGA, an AXI bus
with a router, a small GPU, two compiler toolchains for the core, and
sixteen documents, the tree as it stood on September 16, 2026, were
written in \ganttdays{} days of work spread over four months. This
analysis reconstructs the chronology of those \ganttdays{} days
directly from empirical repository history. Three structural patterns
emerge that diverge from conventional waterfall schedules. The
specification preceded the runtime by four months with no intervening
commits. Implementation proceeded simultaneously across layers rather
than in staged phases, because a project convention required every
concept to enter the runtime, a compiled example, and a document in
the same commit. Foundational infrastructure (including the system bus
and compiler toolchains) arrived last rather than first.
\end{abstract}

\begin{IEEEkeywords}
Software process, project history, hardware description languages,
repository mining.
\end{IEEEkeywords}

\section{What Is Measured}
\label{sec:d-what}

Each band of the chart represents a set of directory paths, and each
block marks a day on which commits touched those paths, with shading
indicating commit density. File paths serve as the empirical metric
because modified paths capture verifiable changes rather than
speculative commit messages. The horizontal axis measures active working
days rather than calendar time, for reasons detailed in
Section~\ref{sec:d-gap}.

The directional dependency arrows are modeled rather than mined.
Because git history records when changes occurred rather than explicit
structural prerequisites, causal dependencies are declared in
\code{//tools/timeline} based on documented architectural dependencies:
the lowering required values and units first, the core required the
lowering, the board required the core, the router required the link,
and the C++ toolchain reused what the Rust toolchain established.

\section{The Four Month Gap}
\label{sec:d-gap}

The project has two lives. \ganttbefore{} commits over
\ganttbeforedays{} days in May are the specification: the language
argued out in prose, with no runtime under it. Then \ganttgapdays{}
days in which nothing was committed at all. Then \ganttafter{}
commits over \ganttafterdays{} days in September, which is
everything else in the tree. The gap is not stated here but found:
\code{//tools/gantt} takes the longest stretch between two working
days and counts what falls on either side of it.

That is why the chart's axis is working days rather than dates. On a
calendar the four idle months would be most of the picture and the
week that holds nine commits in ten would be a smudge a centimetre
wide. The gap is a fact about the project and it is said here in
words, where a reader can weigh it, rather than drawn as an expanse
of nothing.

The gap also explains why the specification band is so thin. It was
finished before the rest began, and it was touched three times
afterwards: twice to publish it as a hermetic article, and once for a
licence header. A document nobody edits is either finished or
forgotten, and which of the two it is cannot be measured from a
repository at all.

\begin{figure*}[t]
\centering
\input{gantt.tex}
\caption{The project by working day. A block is a day on which at
least one commit touched that band's files, shaded by how many did;
the count under each date is the commits of that day across the whole
tree. The arrows are what a band needed before it could begin. The
axis is working days and not the calendar, so the break marks where
it was cut and says what the cut stands for; every other step between
two columns is a night or a weekend. Both the cut's position and its
length are found from the history rather than chosen.}
\label{fig:d-gantt}
\end{figure*}

\section{Everything At Once}
\label{sec:d-together}

In Fig.~\ref{fig:d-gantt}, the activity pattern differs markedly from
a staged waterfall schedule. The runtime, lowering logic, concept
examples, and technical documentation were all active concurrently across
the same four working days.

This concurrency reflects an explicit policy established at project
inception, documented in \code{AGENTS.md} and cited in the
paper~\cite{paper}: every concept must enter the runtime, a compiled
example, and the documentation in the same commit before work is
considered complete. This practice prevented documentation from
drifting from code. Consequently, individual features never appear as
isolated bars within a single band. For example, adding memory lowering
touched \code{lib/macros}, \code{lib/examples}, and \code{docs} in the
same hour, displaying timeline activity across four bands simultaneously.

The bands represent architectural layers rather than chronological
phases, with features cutting vertically across the stack. The chart
confirms that the synchronized development discipline was maintained.

The documentation band remains the widest and densest, with
\ganttdocs{} commits out of \ganttcommits{} across the tree. Beyond the
multi-layer commit rule, this concentration highlights that technical
documents represent primary deliverables rather than secondary
reporting: sixteen documents build continuously from the tree, embedding
source listings directly and generating figures from execution traces.

\section{The Foundations Came Last}
\label{sec:d-last}

The bus, router, GPU, and compiler toolchains appear on the final
working day, reversing the typical sequence where infrastructure is
laid down first.

This ordering was architectural rather than accidental. Each system
relied on components that matured late in development. The AXI link
depended on hardware channel support in \code{Buffer}, which required
conditional send lowering. The router depended on the link; migrating the
core to AXI required the router. Similarly, the C++ toolchain reused the
binary image generator and linker script created for Rust, drastically
reducing secondary bring-up effort.

This progression diverges from conventional top-down design methodology.
The project did not begin by architecting a generalized bus before
implementing a processor. Instead, it built an initial processor around a
minimal dedicated bus, validated it on physical FPGA hardware, and only
subsequently developed the generalized standard bus (decommissioning the
private interconnect thereafter). The processor core served as the
critical vehicle proving language feasibility; standard interconnects
represented a later consolidation step.

\section{What A Repository Cannot Say}
\label{sec:d-cannot}

Three dimensions of effort remain unobservable through repository mining alone.

First, commit counts do not measure wall-clock developer effort. A day
with seventy-one commits reflects different activity than a day with
two, yet neither metric indicates total elapsed problem-solving time.

Second, commit frequency does not distinguish successful paths from
exploratory dead ends. The repository deliberately retains negative
results under \code{experiments/rust\_embedding/}, which appear as standard
commits alongside functional subsystems. Probes with negative outcomes
demanded equivalent exploratory investment.

Third, intra-day commit sequencing is collapsed. Because daily activity
frequently comprised dozens of commits, a daily granularity represents
an intentional abstraction. An hourly axis was evaluated and rejected as
visually cluttered, confirming the necessity of day-level aggregation.

\begin{thebibliography}{9}

\bibitem{paper}
F.~Filmar and D.~Janković, ``TxHDL: A hardware description language
that is a Rust library,'' project repository \code{HDL/txhdl},
\code{//docs:paper}, 2026.

\bibitem{cover}
F.~Filmar and D.~Janković, ``TxHDL documents,'' project repository
\code{HDL/txhdl}, \code{//docs:cover}, 2026.

\end{thebibliography}

\end{document}
